\documentclass[../main.tex]{subfiles}
%DIF LATEXDIFF DIFFERENCE FILE
%DIF DEL /tmp/tmpm_kfke44.tex   Mon Jun 15 17:22:25 2026
%DIF ADD /tmp/tmp6ono493n.tex   Mon Jun 15 17:22:25 2026
%DIF PREAMBLE EXTENSION ADDED BY LATEXDIFF
%DIF UNDERLINE PREAMBLE %DIF PREAMBLE
\RequirePackage[normalem]{ulem} %DIF PREAMBLE
\RequirePackage{color}\definecolor{RED}{rgb}{1,0,0}\definecolor{BLUE}{rgb}{0,0,1} %DIF PREAMBLE
\providecommand{\DIFadd}[1]{{\protect\color{blue}\uwave{#1}}} %DIF PREAMBLE
\providecommand{\DIFdel}[1]{{\protect\color{red}\sout{#1}}} %DIF PREAMBLE
%DIF SAFE PREAMBLE %DIF PREAMBLE
\providecommand{\DIFaddbegin}{} %DIF PREAMBLE
\providecommand{\DIFaddend}{} %DIF PREAMBLE
\providecommand{\DIFdelbegin}{} %DIF PREAMBLE
\providecommand{\DIFdelend}{} %DIF PREAMBLE
\providecommand{\DIFmodbegin}{} %DIF PREAMBLE
\providecommand{\DIFmodend}{} %DIF PREAMBLE
%DIF FLOATSAFE PREAMBLE %DIF PREAMBLE
\providecommand{\DIFaddFL}[1]{\DIFadd{#1}} %DIF PREAMBLE
\providecommand{\DIFdelFL}[1]{\DIFdel{#1}} %DIF PREAMBLE
\providecommand{\DIFaddbeginFL}{} %DIF PREAMBLE
\providecommand{\DIFaddendFL}{} %DIF PREAMBLE
\providecommand{\DIFdelbeginFL}{} %DIF PREAMBLE
\providecommand{\DIFdelendFL}{} %DIF PREAMBLE
%DIF COLORLISTINGS PREAMBLE %DIF PREAMBLE
\RequirePackage{listings} %DIF PREAMBLE
\RequirePackage{color} %DIF PREAMBLE
\lstdefinelanguage{DIFcode}{ %DIF PREAMBLE
%DIF DIFCODE_UNDERLINE %DIF PREAMBLE
  moredelim=[il][\color{red}\sout]{\%DIF\ <\ }, %DIF PREAMBLE
  moredelim=[il][\color{blue}\uwave]{\%DIF\ >\ } %DIF PREAMBLE
} %DIF PREAMBLE
\lstdefinestyle{DIFverbatimstyle}{ %DIF PREAMBLE
	language=DIFcode, %DIF PREAMBLE
	basicstyle=\ttfamily, %DIF PREAMBLE
	columns=fullflexible, %DIF PREAMBLE
	keepspaces=true %DIF PREAMBLE
} %DIF PREAMBLE
\lstnewenvironment{DIFverbatim}[1][]{\lstset{style=DIFverbatimstyle}}{} %DIF PREAMBLE
\lstnewenvironment{DIFverbatim*}[1][]{\lstset{style=DIFverbatimstyle,showspaces=true}}{} %DIF PREAMBLE
\lstset{extendedchars=\true,inputencoding=utf8}

%DIF END PREAMBLE EXTENSION ADDED BY LATEXDIFF

\begin{document}

\begin{center}
    \Large{\textbf{ABSTRACT}}\\
\end{center}
\vspace{1cm}
In the context of the rapidly growing e-commerce sector in Vietnam, small and medium-sized enterprises (SMEs) frequently encounter significant financial and technical barriers when attempting to build and operate independent online stores. Existing e-commerce platforms on the market, such as Shopify, WooCommerce, or Haravan, still exhibit considerable limitations. Specifically, they fail to simultaneously satisfy three crucial criteria: low initial and maintenance costs, flexible interface customization without requiring programming knowledge, and \DIFdelbegin \DIFdel{system scalability capable of serving tens of thousands of stores without linearly increasing infrastructure costs}\DIFdelend \DIFaddbegin \DIFadd{the ability to operate many stores on a shared infrastructure with a low marginal cost for each additional store}\DIFaddend . Stemming from this practical need, this thesis focuses on the research and development of a multi-tenant e-commerce platform named ShopVolo v2 to thoroughly address the resource optimization problem. The primary approach of the study involves implementing a Multi-tenant Software as a Service (SaaS) architecture combined with the "Zero-file" design principle, wherein the entire storefront interface is represented as configuration data rather than discrete physical source code files.

The system is developed based on a Monorepo model using the Turborepo tool, integrating modern technology stacks including Next.js for the administrative dashboard interface and dynamic storefront rendering, and NestJS for the core API system. Furthermore, a dual database system combining PostgreSQL and MongoDB is employed to efficiently process both structured information and flexible configuration documents. The thesis provides three main technical contributions: (i) an automatic multi-tenant data isolation mechanism based on AsyncLocalStorage, ensuring that every query is bound to the correct store scope without passing parameters through processing layers; (ii) a dynamic rendering engine following the Zero-file architecture, allowing new stores to be added without generating any source files or rebuilds; and (iii) an industry-based Master Template model combined with a two-phase draft--publish editing workflow, enabling non-programmers to customize their storefronts safely. The result is a complete system consisting of four applications and five shared library packages, fully implementing the ten proposed functional groups and verified by an automated test suite covering the core business flows.

\end{document}